\subsection{Polynomial Equations in General Forms} 
\begin{procedure}{Solving a Polynomial Equation}
\begin{enumerate*}
\item Write the equation in standard form. This may involve doing multiplication and/or adding or subtracting on both sides.
\item Once the polynomial is in standard form, factor the polynomial side.
\item Use the Zero Property of Multiplication to solve.
\end{enumerate*}
\end{procedure}
\begin{note}
If the polynomial is given in the form $\mbox{factored form}=0$, there is no need to multiply out, then factor again. However, if it is in the form $\mbox{factored form}=c$ for some $c\neq 0$, it \emph{is} necessary to multiply out, subtract $c$ from both sides, then factor again. The eventual factors will be different.
\end{note}
\begin{example}
Solve $(x-3)(x-5)=8$
\begin{lpic}{}{12}
\foreach \x/\step in {-1/1. Write in Standard Form,
											-5/2. Factor:,
											-7/3. Set each factor equal to zero,
											-10/Check each solution}
	\regtextleft{0}{\x}{\step};
\end{lpic}
\end{example}
\begin{example}
Solve $x^3-4x^2+8=x^3-5x^2+3x+6$
\begin{lpic}{}{8}
\foreach \x/\step in {-1/1.,
											-4/2.,
											-5/3.,
											-7/Check:}
	\regtextleft{0}{\x}{\step};
\end{lpic}
\end{example}
